% =============================================================================
% Kapitel 4 -- Implementierung: Relevante Funktionen
% =============================================================================
\section{Implementierung}

Dieses Kapitel beschreibt stichpunktartig die zentralen Funktionen und Subsysteme
von CardSiege. Im Vordergrund steht, was die einzelnen Komponenten leisten, die
zugrundeliegenden Architektur- und Pattern-Entscheidungen wurden in Kapitel~3
erläutert.

% -----------------------------------------------------------------------------
\subsection{Spielzustand und Frame-Loop}
% -----------------------------------------------------------------------------

Der \texttt{GameManager} verwaltet als zentrales Singleton den gesamten
Spielzustand und koordiniert alle Subsysteme. Die Methode \texttt{init(board, cardFactory)}
setzt alle Zustandsvariablen zurück, initialisiert Deck und Hand und startet mit
Startressourcen (3 Energie, 5 Gold, 10 Basislebenspunkte) in der
\texttt{DrawPhase}. Die Frame-Methode \texttt{update(delta)} ruft in fest
definierter Reihenfolge auf: Phasen-Update, Gegnerbewegung,
Reichweitenüberwachung, Turm-Update, Kampfauflösung, Entitätenbereinigung und
Projektilupdates.

Über \texttt{playCard(index)} wird die gewählte Karte aus der Hand entfernt,
Phase und Energieverfügbarkeit geprüft, der Energiebetrag abgezogen und
\texttt{card.execute()} aufgerufen. \texttt{BuildTowerCard}s sind während der
\texttt{WavePhase} gesperrt. Die Methoden \texttt{canPlayCard()} und
\texttt{getEffective\-Energy\-Cost()} kapseln alle Voraussetzungsprüfungen
(Phase, Energie, Gold, vorhandene Türme für Buff-Karten) sowie die Berechnung
aktiver Kartenrabatte.

% -----------------------------------------------------------------------------
\subsection{Turmsystem}
% -----------------------------------------------------------------------------

Über \texttt{queue\-Tower\-Placement(pending)} wird ein \texttt{PendingTower} gespeichert,
sobald ausreichend Gold vorhanden ist; der \texttt{GameScreen} rendert daraufhin
den Platzierungsmarker mit Reichweitenkreis. Beim Klick prüft \texttt{tryPlace\-Pending\-Tower(x, y)}
Baubarkeit des Feldes und ob es bereits belegt ist; bei Erfolg wird
\texttt{PendingTower.build()} aufgerufen, der Turm zur Liste hinzugefügt und Gold
abgezogen.

Beim Upgrade umhüllt \texttt{upgradeTower(tower)} den Turm mit einem
\texttt{UpgradeDecorator} (+1 Schaden, +15 Reichweite, reduzierter Cooldown); die
Kosten steigen pro Level (\texttt{2 + tower.getLevel()} Gold). Globale Buffs wie
\texttt{applyDamageBuff()} ersetzen jeden Turm in der Liste durch einen
entsprechenden Decorator-Wrapper. Für gezielte Buffs auf einzelne Türme speichert
\texttt{queueRangeBuff()} einen Factory-Auftrag, den
\texttt{tryApply\-Pending\-Tower\-Modifier()} beim nächsten vom Spieler ausgewählten
Turm anwendet.

Jeden Frame berechnet \texttt{update\-Range\-Observers()} per Quadratdistanzvergleich,
welche Gegner in Reichweite sind, und benachrichtigt den Turm über
Eintritte und Austritte. In \texttt{resolveCombat(delta)} wird über alle Türme
iteriert:
Bei \texttt{canFire()} wird \texttt{attack\-Strategy.attack()} aufgerufen, der
Cooldown zurückgesetzt und ein visuelles Projektil erzeugt.

% -----------------------------------------------------------------------------
\subsection{Gegnersystem und Kampf}
% -----------------------------------------------------------------------------

In \texttt{Enemy.update(delta)} bewegt sich der Gegner entlang der vorberechneten
Wegpunkt-Liste. Der \texttt{pathIndex} steigt, sobald ein Waypoint erreicht ist;
ein aktiver Slow-Effekt (\texttt{FreezeCard}) skaliert die Geschwindigkeit via
\texttt{slowFactor} für eine begrenzte Dauer. Per Quadratdistanzvergleich prüft \texttt{isInRange(tower)},
ob ein Gegner in Reichweite liegt; der Turmursprung liegt dabei in der
Kachelmitte.

Der \texttt{WaveSpawner} konfiguriert pro Welle drei
\texttt{EnemyPrototype}-Objekte mit skalierten Werten und spawnt Gegner in festen
Zeitintervallen. Sind alle Gegner gespawnt und die Gegnerliste leer, löst
\texttt{endWaveIfComplete()} den Übergang in die \texttt{ResolvePhase} oder
nach zwölf Wellen in die \texttt{GameOverPhase(true)} aus.

Nach jedem Frame bereinigt \texttt{cleanupEntities()} tote Gegner (+1 Gold,
\texttt{EnemyDefeated}-Event) und Gegner, die das Ziel erreicht haben
($-1$ Basislebenspunkt, \texttt{baseDamaged}-Event). Vor dem Entfernen werden
alle Turm-Observer über \texttt{remove\-Enemy\-From\-Observers()} benachrichtigt. Fällt
\texttt{baseHealth} auf null, wechselt das Spiel in \texttt{GameOverPhase(false)}.

% -----------------------------------------------------------------------------
\subsection{Kartensystem}
% -----------------------------------------------------------------------------

Beim Kartenziehen liefert \texttt{Deck.draw()} die oberste Karte; ist der Stapel leer, füllt
\texttt{refill()} ihn über Factory und Prototype-Klonen neu auf und mischt ihn
anschließend. \texttt{Hand} verwaltet die Handkarten als Liste;
\texttt{addToFront()} wird beim Zurücklegen einer abgelehnten Karte genutzt.

Zu Spielbeginn garantiert \texttt{ensure\-First\-Hand\-Has\-Tower()}, dass die erste Hand
mindestens eine \texttt{BuildTowerCard} enthält: Ist keine vorhanden,
wird die erste Handkarte an den Stapelboden gelegt und der Stapel nach einer
Turmkarte durchsucht. Über \texttt{addNext\-Card\-Discount(amount)} werden
Rabatt-Mechaniken der \texttt{DiscountCard} realisiert; der Rabatt wird nach dem nächsten
Kartenspielen auf null zurückgesetzt.

% -----------------------------------------------------------------------------
\subsection{Ereignissystem}
% -----------------------------------------------------------------------------

Der \texttt{GameEventBus} verwaltet eine Liste von
\texttt{Game\-Event\-Listener}-Implementierungen und benachrichtigt alle Subscriber
synchron per \texttt{publish()}. \texttt{GameEvent} ist ein einfaches Wertobjekt
mit typisierten Factory-Methoden (\texttt{towerBuilt(x,y)},
\texttt{phaseChanged(name)}, \texttt{baseDamaged(hp)}, \texttt{EnemyDefeated},
\texttt{cardPlayed(name)}), die Magic-Strings im gesamten System vermeiden. Der
\texttt{GameEventLogger} als einziger aktiver Subscriber schreibt jedes Ereignis
auf die Konsole; weitere Subscriber (Sound, Achievements) könnten ohne Änderung
am Bus per \texttt{subscribe()} ergänzt werden.

% -----------------------------------------------------------------------------
\subsection{Spielfeld und Pfadgenerierung}
% -----------------------------------------------------------------------------

\texttt{Tile} repräsentiert eine einzelne Kachel mit Gitterkoordinaten und
einem booleschen \texttt{path}-Flag. \texttt{Board} hält die Kachelliste,
die vorberechnete Pfad-Vektorliste und stellt \texttt{isBuildable(gridX,~gridY)}
für Bauprüfungen bereit.

Der Pfad wird einmalig in \texttt{Board.createDefault()} erzeugt: Zwölf
relativ zur Spielfeldgröße skalierte Kontrollpunkte werden durch
ausschließlich horizontale und vertikale Segmente verbunden. Die resultierenden
Gitterkoordinaten landen als \texttt{GridPoint}-Records in einem
\texttt{LinkedHashSet}, das Duplikate entfernt und die Reihenfolge erhält. Jede Koordinate wird als
\texttt{Tile(path=true)} markiert und als Weltkoordinate in die Pfad-Liste
übernommen, entlang derer Gegner sich bewegen.

% -----------------------------------------------------------------------------
\subsection{Projektile}
% -----------------------------------------------------------------------------

Als reines Darstellungsobjekt speichert \texttt{Projectile} Start- und Zielposition
sowie eine \texttt{ttl}-Zeit von 0,15~Sekunden; die Spiellogik bleibt davon
unberührt. Für weiches Ausblenden liefert \texttt{getAlpha()} einen linear
abfallenden Transparenzwert; \texttt{isExpired()}
signalisiert dem \texttt{GameManager} das Entfernen aus der Liste. Schaden und
Trefferauswertung sind bereits in \texttt{resolveCombat()} abgeschlossen,
bevor das Projektil erzeugt wird.

% -----------------------------------------------------------------------------
\subsection{Benutzeroberfläche und Eingabe}
% -----------------------------------------------------------------------------

Der \texttt{GameScreen} bildet die einzige Verbindung zwischen Spiellogik und
Darstellung. Pro Frame delegiert \texttt{render()} zunächst \texttt{update(delta)} an
den \texttt{GameManager} und zeichnet danach Hintergrund, Spielfeld, Entitäten,
Platzierungsvorschau und HUD. Türme erscheinen als farbige Rechtecke, Gegner
als Kreise mit HP-Balken, Projektile als gelbe Linien mit \texttt{alpha}-Ausblenden.

Aktiv ist \texttt{drawPlacementPreview()}, solange ein \texttt{PendingTower}
im Manager registriert ist: Das überfahrene Feld wird grün (baubar) oder rot
(gesperrt) eingefärbt, ein Reichweitenkreis zeigt die Abdeckung -- ohne dass
ein echter Turm existiert. Das HUD zeigt Ressourcen und Handkarten; spielbare
Karten werden grün, nicht spielbare rot hinterlegt. Über \texttt{collectBuffs()}
wird die Decorator-Kette eines selektierten Turms traversiert und alle
aktiven Boni für das Tooltip summiert.

Die innere Klasse \texttt{GameInputProcessor} übernimmt die Eingabe:
\texttt{[1--8]} spielen Karten, \texttt{SPACE} startet die Welle, \texttt{U}
upgraded den selektierten Turm, \texttt{R} startet neu und \texttt{ESC} kehrt
zum Menü zurück. Mausklicks werden per \texttt{camera.unproject()} in
Gitterkoordinaten konvertiert und lösen je nach Zustand
\texttt{tryPlace\-PendingTower()} oder \texttt{tryApply\-Pending\-Tower\-Modifier()} aus.
